\documentclass[UTF8,11pt,a4paper,fontset=fandol]{ctexart}

\usepackage[margin=2.35cm]{geometry}
\usepackage{amsmath,amssymb,mathtools}
\usepackage{booktabs}
\usepackage{tabularx}
\usepackage{array}
\usepackage{enumitem}
\usepackage{xcolor}
\usepackage{tcolorbox}
\usepackage{fancyhdr}
\usepackage{hyperref}
\usepackage{caption}

\definecolor{Ink}{HTML}{172033}
\definecolor{Muted}{HTML}{5B6475}
\definecolor{DeepBlue}{HTML}{0B3D91}
\definecolor{SoftBlue}{HTML}{EEF5FF}
\definecolor{SoftGreen}{HTML}{EFFAF3}
\definecolor{SoftAmber}{HTML}{FFF7E6}
\definecolor{SoftRose}{HTML}{FFF0F3}
\definecolor{RuleGray}{HTML}{D8DEE9}

\hypersetup{
  colorlinks=true,
  linkcolor=DeepBlue,
  urlcolor=DeepBlue,
  citecolor=DeepBlue
}

\pagestyle{fancy}
\fancyhf{}
\setlength{\headheight}{24pt}
\lhead{\textcolor{Muted}{总结}}
\rhead{\textcolor{Muted}{\thepage}}
\renewcommand{\headrulewidth}{0.4pt}
\renewcommand{\headrule}{\hbox to\headwidth{\color{RuleGray}\leaders\hrule height \headrulewidth\hfill}}

\setlist[itemize]{leftmargin=1.8em,itemsep=0.28em,topsep=0.35em}
\setlist[enumerate]{leftmargin=2em,itemsep=0.35em,topsep=0.35em}
\captionsetup{font=small,labelfont=bf}

\newtcolorbox{keybox}[1]{
  colback=SoftBlue,
  colframe=DeepBlue,
  boxrule=0.6pt,
  arc=2mm,
  left=1.1em,
  right=1.1em,
  top=0.8em,
  bottom=0.8em,
  title=\textbf{#1}
}

\newtcolorbox{decisionbox}[1]{
  colback=SoftGreen,
  colframe=green!45!black,
  boxrule=0.55pt,
  arc=2mm,
  left=1.1em,
  right=1.1em,
  top=0.75em,
  bottom=0.75em,
  title=\textbf{#1}
}

\newtcolorbox{riskbox}[1]{
  colback=SoftAmber,
  colframe=orange!65!black,
  boxrule=0.55pt,
  arc=2mm,
  left=1.1em,
  right=1.1em,
  top=0.75em,
  bottom=0.75em,
  title=\textbf{#1}
}

\title{\vspace{-1.8em}\textbf{\Huge 多智能体规模、拓扑与合并负担}\\[0.35em]
\Large 从局部 Error Rate 到图结构性能函数的实验路线}
\author{Agent Expressional Graph}
\date{\today}

\begin{document}
\maketitle

\begin{keybox}{一句话总结}
本次形成的核心问题是：在 Count Frequency 等可拆分任务中，
增加 agent 数量会降低单个 agent 面对的局部问题规模，
从而可能降低局部 error rate；但同时也会增加跨 agent 合并、对齐和聚合的负担。
接下来的实验目标，是寻找 agent 数量、通信 topology 与 merge burden 之间的平衡点，
并尝试建立从图结构到最终 error 的函数关系。
\end{keybox}

\section{背景与动机}

当前讨论围绕一个直观但关键的 trade-off 展开。
对于一个固定规模的任务，如果使用更多 agents，每个 agent 分到的子问题会变小。

例如在 Count Frequency 问题中，数组被切成更多 subarray 后，每个 agent 只需要处理更少的数字。
这样理论上会降低单个 agent 内部出错的概率，因为局部上下文更短、局部计数或分类判断更简单。

但这个收益并不是免费的。
agent 数量增加后，系统末端需要合并更多局部结果，
merge 的复杂度、通信量和错误传播路径都会增加。

因此，问题不只是“agent 越多越好吗”，而是：

\[
\textbf{局部问题规模下降带来的收益}
\quad \text{是否能够超过} \quad
\textbf{全局合并负担上升带来的损失}.
\]

\section{核心问题}

\begin{decisionbox}{核心研究问题}
是否存在一个理论平衡的 agent 数量 \(N^\star\)，
使得局部求解误差和最终 merge burden 之间达到较优折中？
进一步地，这个平衡点是否依赖于通信 topology？
\end{decisionbox}

可以把最终系统误差拆成两个来源：

\[
e_{\mathrm{final}}
=
e_{\mathrm{local}}
+
e_{\mathrm{merge}}
+
e_{\mathrm{propagation}},
\]

其中：

\begin{itemize}
  \item \(e_{\mathrm{local}}\)：agent 在自己局部数组或局部子任务上的错误；
  \item \(e_{\mathrm{merge}}\)：多个局部结果在合并时产生的遗漏、重复计数、冲突或格式不一致；
  \item \(e_{\mathrm{propagation}}\)：错误信息在 topology 上传播后被其他 agents 接受或放大的部分。
\end{itemize}

增加 agent 数量通常会降低 \(e_{\mathrm{local}}\)，但可能提高 \(e_{\mathrm{merge}}\) 和 \(e_{\mathrm{propagation}}\)。因此，一个可操作的实验目标是寻找：

\[
N^\star
=
\arg\min_N
\left[
e_{\mathrm{local}}(N)
+
e_{\mathrm{merge}}(N,\tau)
+
e_{\mathrm{propagation}}(N,\tau)
\right],
\]

其中 \(\tau\) 表示通信 topology。

\section{DIG 视角：哪些 Agents 真正有用}

提出了一个自然切入点：给定一组 agents 运行任务后，不只看最终答案，还要观察系统形成的交互图结构。
关键问题是：

\begin{quote}
在一个 MAS run 中，真正对最终答案有贡献的 agents 有几个？这些 agents 在图中处于什么位置？
\end{quote}

这实际上接近 DIG（Dynamic Interaction Graph）的思想：
把每轮交互、消息传递、合并和最终贡献都记录为图上的结构信号。这样可以从图中抽取特征，例如：

\begin{itemize}
  \item 有效参与的 agent 数量；
  \item 最终答案依赖的局部结果数量；
  \item merge 节点的 fan-in；
  \item 信息传播路径长度；
  \item 是否存在少数高贡献 agents；
  \item 错误是否从某些局部节点扩散到全局结果。
\end{itemize}

\begin{keybox}{图结构问题}
如果把 topology 或运行后形成的交互图记为 \(E\)，把最终 error 记为 \(e\)，我们希望收集大量 pair：
\[
(E,e),
\]
并学习或拟合一个从图结构到 error 的函数：
\[
e=f(E).
\]
进一步地，也可以把问题反过来：如果已经观测到某种 error pattern，例如局部 error 很低但最终 error 很高，或者少数 agents 的错误被全局放大，那么能否近似构造一个逆向映射：
\[
E \approx f^{-1}(e),
\]

\end{keybox}

\section{从图到 Error 的建模路径}

本次讨论中，最困难但也最有研究价值的问题是如何建立从图结构到 error 的函数。
直接学习 \(e=f(E)\) 比较难。
GNN可以但是需要尝试拓展GNN之外的方法。

\section{下一阶段实验 Setup}

下一步重点观察 Count Frequency（CF）问题。初始设置如下：

\begin{itemize}
  \item 全局数组规模：\(M=5000\)；
  \item 数字范围：\(1\) 到 \(1000\)；
  \item agents 只看到分配给自己的 subarray；
  \item 尝试不同 topology；
  \item 每轮进行一次 merge 或 aggregation；
  \item 每轮记录总体 error rate 和每个 agent 的局部 error rate；
  \item setup 参考 MANCET 论文中的多智能体评测范式：固定任务族，控制 agent 数量、轮数和 topology，逐轮记录系统状态。
\end{itemize}

\begin{table}[h]
\centering
\caption{实验控制变量}
\begin{tabular}{lll}
\toprule
变量 & 取值 & 目的 \\
\midrule
Agent 数量 \(N\) & \(2,4,8,16,\ldots\) & 观察规模与合并负担的 trade-off \\
通信轮数 \(R\) & \(2,3,5\) & 观察短程通信是否足够降低误差 \\
数组长度 \(M\) & \(5000\) & 固定任务总规模 \\
数字范围 & \(1\) 到 \(1000\) & 固定类别/值域规模 \\
Topology \(\tau\) & chain, star, mesh, exponential 等 & 比较不同通信结构 \\
Merge 频率 & 每轮一次 & 观察逐轮合并带来的收益和损失 \\
\bottomrule
\end{tabular}
\end{table}

\end{document}
